\documentclass[11pt]{article}
\usepackage[a4paper,margin=25mm]{geometry}
\usepackage{amsmath,amssymb,amsthm,booktabs}
\usepackage[colorlinks=true,urlcolor=blue,citecolor=blue]{hyperref}
\newtheorem{theorem}{Theorem}
\newtheorem{lemma}{Lemma}
\newcommand{\C}{\mathcal C}
\title{A sharp threshold for three-colour targets from four-colour cubes}
\author{Jason Davies}
\date{November 1, 2026}
\begin{document}
\maketitle

\begin{abstract}
Any collection of 29 cubes, each displaying all four colours of a fixed
palette, contains eight that form a $2\times2\times2$ cube with monochromatic
faces using exactly three colours. There is a collection of 28 cubes with
no such solution. This answers the threshold question in Problem 2 of Berkove, Eickemeyer and Schulman
(2018), subject to a further priority check. Moreover, there are exactly two
labelled 28-cube obstructions, exchanged by reflection. The proof uses their reduction
by corner containment and elementary counting. The threshold and classification are formalised in Lean.
\end{abstract}

\noindent\textbf{AI transparency statement.} The research argument, computer
code, Lean formalisation and manuscript were developed using AI under the
author's direction. The declaration at the end gives the scope of this work
and its review status.

\section{Model and geometric facts}
Colours are fixed, varieties are identified by proper rotations, and
reflections are not allowed as moves. Collections are multisets: repeated
varieties are allowed. Only exterior faces must match the target; internal
faces are unrestricted. A corner is a cyclically ordered triple, so $abc$
and $acb$ are distinct when the three colours differ. A cube supplies one
corner position, regardless of how many suitable corners it has.

Write $Q(a;b,c)$ for a target whose $b$ and $c$ faces are opposite and whose
other four faces have colour $a$. Its corner multiset is
$4aab+4aac$. Write $R(a;b,c)$ for the analogous target with $b$ and $c$
adjacent; its corners are
\[
  2aaa+2aab+2aac+abc+acb.
\]
Finally, a target with opposite pairs $aa,bb,cc$ has four corners of each
orientation on $\{a,b,c\}$. All three targets use exactly three colours.

If every corner type of a variety $v$ occurs on a variety $w$, replacing a
copy of $w$ by $v$ cannot create a solution that was not possible before:
each small cube is used at just one corner. This is the reduction of
\cite[Section 2]{BES}. For clarity, Table~\ref{tab:reduction} gives the
entire reduction needed here. A face tuple is read in the order
$+x,-x,+y,-y,+z,-z$. The rows represent the seven orbits under colour
permutations, with rotations already identified. The three non-retained
representatives all contain the corner types of $001123$. Colour
permutations transport these containments to every member of their orbits.

\begin{table}[ht]
\centering
\begin{tabular}{ccc}
\toprule
Representative & Labelled varieties in orbit & Replacement or retained type\\
\midrule
$000123$ & 12 & $001123$\\
$001123$ & 6  & checkerboard $C$\\
$001213$ & 12 & $001123$\\
$010123$ & 6  & $001123$\\
$010203$ & 8  & triple-corner $T$\\
$010213$ & 12 & pair type $P$\\
$010223$ & 12 & pair type $P$\\
\bottomrule
\end{tabular}
\caption{The 68 varieties reduce to 38. The orbits can be obtained by
placing the repeated colours: either a triple, or two pairs.}
\label{tab:reduction}
\end{table}

The retained types are precisely the minimal varieties in
\cite[Table 8]{BES}. We use the following explicit description.
\begin{itemize}
\item The eight $T$ varieties have a monochromatic corner $aaa$, and also
  $aab$ for each of the other three colours $b$. Each belongs to the
  bucket of its repeated colour $a$.
\item The 24 $P$ varieties belong to two buckets, one for each repeated
  colour. In bucket $a$, such a cube has exactly two of the three types
  $aab$ with $b\ne a$.
\item The six checkerboard varieties have two repeated colours on
  opposite pairs of faces and two singleton colours on the remaining
  opposite pair. A checkerboard whose repeated colours are $a,b$ has both
  orientations on the triples $\{a,b,c\}$ and $\{a,b,d\}$. It belongs to
  no bucket.
\end{itemize}
The 32 ordinary varieties ($T$ and $P$) split into two classes of 16.
Every member of one class has the same four trichromatic corner types,
one orientation for each three-element subset of the palette; every
member of the other class has the four reversed orientations.
Each class contains exactly one $T$ variety for each repeated colour.
For our
coordinate convention the first four types can be taken as
\[
 \Omega=\{012,031,023,132\}.
\]
These properties follow by reading the corners of the four retained
representatives and permuting colours. A separate geometric checker
verifies the table and all the stated incidences using physical rotations.

\begin{lemma}[Bucket bound]\label{lem:bucket}
Eight reduced cubes in a single bucket can build a three-colour target.
Consequently, a reduced collection with no such solution has at most seven
cubes in each bucket.
\end{lemma}
\begin{proof}
For bucket $a$, count the cubes supporting each of the three types $aab$.
Each cube supports at least two types, giving at least 16 incidences.
At least two types, say $aab$ and $aac$, each have support at least four:
otherwise the total would be at most $8+3+3=14$. Every cube supports at
least one of these two types. Thus eight cubes can be matched to four
$aab$ and four $aac$ positions, constructing $Q(a;b,c)$. The matching
conditions are just support at least four for each type and at least
eight for their union.
\end{proof}

\section{The extremal collections}
\begin{theorem}\label{thm:classification}
A collection of 28 four-colour cubes without a three-colour corner
solution consists of seven copies of each of the four $T$ varieties
in one ordinary orientation class.
\end{theorem}
\begin{proof}
Reduce the collection as above and suppose that no solution exists.
Let $T,P,C$ be the numbers of triple-corner, pair-type and checkerboard
cubes. Let $A\ge B$ be the sizes of the two ordinary orientation classes,
so $A+B=T+P$ and $T+P+C=28$.

For a three-element colour set $E$, let $z_E$ count checkerboards supporting
its two orientations. If $z_E\ge8$, those eight checkerboards build the
target with three monochromatic opposite pairs. Thus $z_E\le7$ for each
$E$, and, since each checkerboard supports two triples, $2C\le28$.
It follows that $A+B\ge14$, so $A\ge7$. If $B+z_E\ge4$, four cubes
from the smaller class and the eligible checkerboards supply one
orientation, and four unused cubes from the larger class supply the other.
Hence $B+z_E\le3$ for every $E$. Summing gives
\begin{equation}\label{eq:checker}
 C+2B\le6,\qquad C\le6,\qquad A=28-C-B\ge22.
\end{equation}

Let $t_a$ and $p_a$ count $T$ and $P$ cubes in bucket $a$. By
Lemma~\ref{lem:bucket}, $t_a+p_a\le7$. Summing over the four colours gives
\begin{equation}\label{eq:buckets}
 T+2P\le28,\qquad C\ge P.
\end{equation}
Suppose first that a checkerboard is present. For every colour $a$ it supports
some triple $\{a,b,c\}$. If $t_a\ge6$, six $T$ cubes supply
$2aaa+2aab+2aac$, the checkerboard supplies one trichromatic orientation,
and an unused cube of the larger ordinary class supplies the other.
There is such an unused cube because $A\ge22$. This constructs
$R(a;b,c)$, a contradiction. Therefore
\begin{equation}\label{eq:triple}
 t_a\le5,\quad T\le20,\quad P+C=28-T\ge8.
\end{equation}
Together with \eqref{eq:checker} and \eqref{eq:buckets}, this implies
\begin{equation}\label{eq:range}
 4\le C\le6.
\end{equation}

We claim that $t_a\ge4$ forces $t_a+p_a\le5$.
For any two colours $a,b$, some checkerboard supports a triple containing
both: otherwise every checkerboard would have the complementary pair of
repeated colours, and one of the corresponding $z_E$ would equal
$C\ge4$, contradicting $z_E\le3$.
If $t_a=5$ and $p_a\ge1$, choose an $aab$ corner on a pair-type cube and
a checkerboard supporting $\{a,b,c\}$. The pair-type cube and five $T$
cubes supply the six non-trichromatic corners of $R(a;b,c)$.
If $t_a=4$ and $p_a\ge2$, two pair-type cubes have a common third colour
$b$, since each supports two of the three possibilities. These two cubes
supply $2aab$, and four $T$ cubes supply $2aaa+2aac$, where a checkerboard
again supports $\{a,b,c\}$. In either case that checkerboard and an unused
cube of the larger ordinary class fill the last two corners. This proves
the claim.

For $t_a\le3$, the bucket bound gives $t_a+p_a\le7\le10-t_a$.
For $4\le t_a\le5$, the claim gives $t_a+p_a\le5\le10-t_a$.
Summing over all four colours therefore gives
\[
 T+2P\le40-T,\qquad T+P\le20.
\]
But $T+P=28-C\ge22$, a contradiction. Thus $C=0$, and
\eqref{eq:buckets} then forces $P=0$. All cubes are of type $T$,
and each $t_a=7$.

If both ordinary orientation classes occur, take a cube from the smaller
class, with repeated colour $b$, and choose $a\ne b$. Six $T$ cubes
of repeated colour $a$ supply the six non-trichromatic positions of
$R(a;c,d)$, for any distinct $c,d\ne a$. The chosen smaller-class cube
and an unused larger-class cube supply the two remaining orientations.
This is again a solution. All 28 cubes must therefore belong to one class.

Finally, this conclusion holds for the original collection, not just its
reduction. Every non-retained variety in Table~\ref{tab:reduction} is
replaced by a checkerboard; since the reduced collection contains none,
no such replacement could have occurred.
\end{proof}

\begin{theorem}\label{thm:upper}
Every collection of 29 four-colour cubes contains a three-colour corner
solution.
\end{theorem}
\begin{proof}
Otherwise every 28-cube subcollection would have the form in
Theorem~\ref{thm:classification}. Take one such subcollection $S$ and let
$v$ be the remaining cube. Remove from $S$ a cube of a variety different
from $v$, and insert $v$. The resulting 28 cubes are not seven copies of
each of the four varieties in the original class. They cannot belong
entirely to the other class either, since 27 cubes of $S$ remain. This
contradicts Theorem~\ref{thm:classification}.
\end{proof}

\section{Sharpness}
\begin{lemma}\label{lem:balance}
In a cube using exactly three colours, the two orientations of a
trichromatic corner occur equally often. If no trichromatic corner occurs,
the cube is $Q(a;b,c)$ for some distinct $a,b,c$.
\end{lemma}
\begin{proof}
Colour the vertices of the dual octahedron by their corresponding face
colours and orient its triangular faces consistently. Map every directed
edge between two distinct colours to the corresponding directed edge of
an abstract triangle on $a,b,c$, and an equal-colour edge to zero.
The sum of the triangular boundaries is zero, since each octahedron edge
occurs twice with opposite orientations. A triangle with a repeated
colour contributes zero; a trichromatic triangle contributes either
$[a,b]+[b,c]+[c,a]$ or its negative. Their numbers are therefore equal.

For the second assertion, if a colour $c$ occurs on only one face, its
four neighbours must all have the same colour $a$: a change from $a$ to
$b$ around this face would produce a trichromatic corner. The opposite
face is then the sole $b$ face. Otherwise each colour occurs twice.
If the two $c$ faces are opposite, the four neighbours of either contain
both $a$ and $b$, forcing a change. If they are adjacent, the other three
neighbours form a path containing both $a$ and $b$, again forcing a
change. Both cases produce a trichromatic corner.
\end{proof}

\begin{theorem}\label{thm:sharp}
The threshold 29 in Theorem~\ref{thm:upper} is best possible. Exactly two
labelled 28-cube collections fail, and they are mirror images.
\end{theorem}
\begin{proof}
Take seven copies of each of the following four face tuples:
\[
 010203,\qquad 011213,\qquad 021223,\qquad 031323.
\]
These are $T$ cubes with repeated colours $0,1,2,3$, respectively, all in
the same ordinary orientation class. A target with any trichromatic corner
requires both orientations by Lemma~\ref{lem:balance}, whereas the
collection supplies only one. A target without one is $Q(a;b,c)$;
every one of its eight corners has colour $a$ twice, so only the seven
$T$ cubes of repeated colour $a$ are usable. No three-colour target is
possible. Reflection gives the analogous example in the other orientation
class. Theorem~\ref{thm:classification} excludes all other collections.
\end{proof}

\section{Formal verification and computation}
The threshold and the full classification have been formalised in Lean
4.35.0~\cite{Lean4} with mathlib v4.35.0~\cite{mathlib}. The formal model assigns arbitrary natural-number
multiplicities to the 68 four-colour varieties. A physical solution selects
eight distinct copies and specifies a proper rotation matching the three
exposed faces at each corner of a target. Separate kernel checks establish
coverage of all four-colour source and three-colour target face assignments,
uniqueness of source varieties up to rotation, and the equivalence of corner
compatibility with physical face matching.

The formal upper bound follows the counting argument, using finite Hall
certificates for the sufficient constructions. The lower bound checks a Hall
deficiency for each of the 120 targets and each extremal collection;
reflection exchanges the two collections. Thus the formal proof does not
enumerate all collections. Its axiom audit reports only Lean's standard
\texttt{propext}, \texttt{Classical.choice}, and \texttt{Quot.sound}; it uses
neither proof holes nor \texttt{native\_decide}. The accompanying repository\footnote{\url{https://github.com/jasondavies/macmahon-four-to-three}}
contains the source and reproduction commands. It reuses the proper-rotation
and outward-face geometry from \href{https://github.com/jasondavies/macmahon-min-universal-sets}{the earlier MacMahon formalisation}, pinned to
commit \texttt{3f854795a77cd7c902e4c5a1a5da3322f9cd44fa}. No classification
result from that project is used.

The proof does not depend on optimisation. A Hall-obstruction model first
found the 28-cube example. After the
bucket reduction, the compiled multicore CP-SAT backend also reported
the maximum obstruction size as 28. That solver status is a computational check, not a separately
verified optimality certificate.

Two further feasibility runs reported infeasibility for a 29-cube
obstruction and for a 28-cube obstruction other than the two specified
collections. Both use the reduced source family and the bucket bounds.

The accompanying script \texttt{check\_three\_colour\_threshold.py}
independently generates the 24 rotations from quarter-turns, verifies the
38-type reduction and all corner incidences, and checks the lower example
against all 120 three-colour targets on the four-colour palette by
bipartite matching and explicit Hall deficiencies. It also checks all
68 one-cube extensions of the lower example and 1,000 sampled collections
of size 29. These samples do not establish the upper bound; the written
proof does.

Problem 2 of \cite{BES} explicitly asks for this threshold. An expanded,
bounded literature search on November 1, 2026 found no later answer, but does
not establish priority. In particular, Rosenfeld's October 2026 article
\cite{Rosenfeld} was identified but could not be read. A separate AI review
found no mathematical errors in the argument or physical model. These
automated checks are not independent human peer review.

\section*{Declaration of AI-generated research and review}
GPT-6-Astra (OpenAI), used through Codex, developed the research argument,
computational search, Lean formalisation and manuscript under the author's
direction. A separate AI reviewer checked the argument and physical model.
The human author initiated and directed the research and accepts
responsibility for the content of this paper and the accuracy of this
disclosure. The paper has not received
independent human peer review; formal verification does not establish
publication priority.

\begin{thebibliography}{9}
\raggedright
\bibitem{BES}
E. Berkove, A. Eickemeyer and M. Schulman,
\emph{The 4-Color Cubes Puzzle}, Integers \textbf{18} (2018), A68.
\href{https://math.colgate.edu/~integers/s68/s68.pdf}{Journal PDF}.
\bibitem{Rosenfeld}
M. Rosenfeld, \emph{MacMahon Cubes and Moshe's Insanity, Puzzle 1},
Geombinatorics Quarterly \textbf{36}(2) (October 2026), p.~78 ff.
\href{https://geombina.uccs.edu/past-issues/volume-xxxvi}{Publisher listing}
(full text not examined).
\bibitem{Lean4} L.~de Moura and S.~Ullrich,
\emph{The Lean 4 Theorem Prover and Programming Language},
in Automated Deduction---CADE~28, LNCS \textbf{12699}, Springer (2021),
625--635.
\href{https://doi.org/10.1007/978-3-030-79876-5_37}{doi:\nolinkurl{10.1007/978-3-030-79876-5_37}}.
\bibitem{mathlib} The mathlib Community,
\emph{The Lean Mathematical Library},
in Proceedings of the 9th ACM SIGPLAN International Conference on
Certified Programs and Proofs (CPP~2020), ACM (2020), 367--381.
\href{https://doi.org/10.1145/3372885.3373824}{doi:\nolinkurl{10.1145/3372885.3373824}}.
\href{https://arxiv.org/abs/1910.09336}{arXiv:1910.09336}.
\end{thebibliography}
\end{document}
